

\documentclass[11pt,a4paper]{article}
\usepackage[utf8]{inputenc}
\usepackage[T1]{fontenc}
\usepackage[slovak]{babel}
\usepackage{geometry}
\usepackage{tabularx}
\usepackage{longtable}
\usepackage{array}
\usepackage[table]{xcolor}
\usepackage{booktabs}
\usepackage{hyperref}
\usepackage{enumitem}
\geometry{margin=2.2cm}
\sloppy
\setlength{\emergencystretch}{3em}
\renewcommand\tabularxcolumn[1]{p{#1}}
\setlist[itemize]{leftmargin=1.2em, topsep=3pt, itemsep=1pt}
\setlist[enumerate]{leftmargin=1.5em, topsep=3pt, itemsep=1pt}

\definecolor{apvvgray}{RGB}{240,240,240}
\newcolumntype{Y}{>{\raggedright\arraybackslash}X}
\newcommand{\Field}[2]{\rowcolor{apvvgray}\textbf{#1} & #2\\}


\begin{document}
\begin{center}
{\Large \textbf{Žiadosť pre aplikovaný výskum a vývoj}}\\[2pt]
{\large \textbf{Formulár VV–A až VV–E (Príloha 5)}}\\[4pt]
\textit{Projekt: TIMA – TimeManager: Inteligentný plánovač času pre študentov s integráciou Google služieb}
\end{center}

% ===================== A1 =====================
\section*{VV–A1: Základné informácie o projekte}
\begin{tabularx}{\textwidth}{@{}lY@{}}
\Field{Evidenčné číslo projektu}{APVV-24-XXXX (pridelí agentúra)}
\Field{Dátum podania}{doplní žiadateľ}
\Field{Názov projektu}{TimeManager: Inteligentný plánovač času pre študentov s integráciou Google služieb (TIMA)}
\Field{Akronym}{TIMA}
\Field{Odbor vedy a techniky}{9.2 Informatika a informačné technológie – Softvérové inžinierstvo, umelá inteligencia}
\Field{Charakter výskumu}{Aplikovaný výskum (vývoj softvéru a validácia v relevantnom prostredí)}
\Field{Začiatok riešenia}{01.01.2026}
\Field{Koniec riešenia}{31.12.2028}
\Field{Anotácia (max. 2\,000 znakov)}{Projekt TIMA vyvíja a overuje inteligentnú webovú aplikáciu na plánovanie času pre študentov. Riešenie kombinuje osvedčené metódy time-managementu (Eisenhowerova matica, Pomodoro technika) s integráciou Google Calendar/Tasks a AI (Gemini) pre kontextové odporúčania. Architektúra: frontend Next.js 14 (App Router, RSC, Tailwind, shadcn/ui), backend Django 5 + DRF, databáza PostgreSQL 16 (JSONB, full-text). Kľúčové prínosy: inteligentná prioritizácia a obojsmerná synchronizácia, vizuálne dashboardy, Pomodoro timer, tímová spolupráca. Projekt cieli na pilotnú prevádzku na univerzitnom prostredí a merateľné zlepšenie produktivity študentov (čas do dokončenia úlohy, adherence na plán, subjektívna spokojnosť).}
\Field{Žiadateľská organizácia}{Slovenská technická univerzita v Bratislave, Fakulta informatiky a informačných technológií (FIIT STU)}
\end{tabularx}

\vspace{0.6em}
\begin{tabularx}{\textwidth}{@{}lY@{}}
\Field{Požadované finančné prostriedky z APVV (EUR)}{239\,360}
\Field{Spolufinancovanie projektu (EUR)}{42\,240}
\Field{Celkové náklady na projekt (EUR)}{281\,600}
\Field{Poznámka}{Rozpočet je orientačný a bude spresnený podľa kalkulácie ľudských zdrojov a služieb.}
\end{tabularx}

% ===================== A2 =====================
\section*{VV–A2: Základné informácie o riešiteľských organizáciách}
\textbf{Žiadateľ – Applicant}

\begin{tabularx}{\textwidth}{@{}lY@{}}
\Field{Názov organizácie}{Slovenská technická univerzita v Bratislave, FIIT STU}
\Field{Skrátený názov}{FIIT STU}
\Field{Adresa}{Ilkovičova 2, 842 16 Bratislava, Slovenská republika}
\Field{IČO}{00397687}
\Field{Príslušnosť k rezortu}{Vysoké školy}
\Field{Forma hospodárenia}{Verejná vysoká škola}
\Field{Kontaktná osoba / E-mail}{Michal Pípa / michal.pipa@stuba.sk}
\Field{Telefón}{+421 905 000 000}
\end{tabularx}

\smallskip
\noindent \textbf{Spoluriešiteľská organizácia:} v tomto návrhu neplánovaná (možno doplniť neskôr).

% ===================== A3 =====================
\section*{VV–A3: Zoznam riešiteľov / List of participants}
{\small
\begin{tabularx}{\textwidth}{@{}p{3.6cm}p{1.7cm}Xp{2.0cm}p{2.0cm}p{1.7cm}p{1.7cm}@{}}
\rowcolor{apvvgray}\textbf{Meno a priezvisko} & \textbf{Tituly} & \textbf{Pracovné zaradenie} & \textbf{Dátum nar.} & \textbf{IČO org.} & \textbf{Hodín} & \textbf{Hod./roky}\\ \midrule
Michal Pípa & Bc. & Zodpovedný riešiteľ / vývoj backendu a AI &  & 00397687 & 1800 & 0.9\\
Michal Urban &  & Frontend vývoj &  & 00397687 & 1500 & 0.75\\
Tomáš Topalov &  & Integrácie Google API &  & 00397687 & 1000 & 0.5\\
Tomáš Šuba &  & Výskum UX a evaluácia &  & 00397687 & 1000 & 0.5\\
\end{tabularx}
}
% ===================== A4 =====================
\section*{VV–A4: Zodpovedný riešiteľ}
\begin{tabularx}{\textwidth}{@{}lY@{}}
\Field{Meno a priezvisko}{Michal Pípa}
\Field{Pohlavie}{—}
\Field{Telefón}{+421 905 000 000}
\Field{E-mail}{michal.pipa@stuba.sk}
\Field{Zodpovedný riešiteľ je mladým vedeckým pracovníkom}{Áno}
\Field{Vedecké databázy / ORCID / WoS / Scopus}{doplní žiadateľ}
\end{tabularx}

% ===================== B =====================
\section*{VV–B: Ciele, zámery a výstupy projektu}
\textbf{Kľúčové slová:} time management, Eisenhowerova matica, Pomodoro, Google Calendar, Google Tasks, Gemini, Next.js, Django, PostgreSQL, študenti.

\textbf{Ciele projektu:}
\begin{enumerate}
\item Implementovať backend (Django + DRF) a dátový model do 12 mesiacov.
\item Vyvinúť AI modul (Gemini) pre prioritizáciu, odhad času a plánovanie.
\item Dodať frontend (Next.js 14, Tailwind, shadcn/ui) s dashboardom a Pomodoro.
\item Integrovať Google Calendar/Tasks s obojsmernou synchronizáciou.
\item Realizovať pilot a evaluáciu (min. 200 používateľov) s reportom dopadu.
\end{enumerate}

\noindent\textbf{TRL projektu:} TRL 5–6.\\
\textbf{Zdôvodnenie TRL:} Funkčný prototyp integrovaný s reálnymi službami (Google API) bude overený v relevantnom prostredí na univerzite. Napĺňa sa validácia subsystémov a prototypu (TRL5) a následne demonštrácia v relevantnom prevádzkovom prostredí s pilotnými používateľmi (TRL6).

\noindent\textbf{Verejné výstupy podľa etáp:}
\begin{itemize}
\item WP1 (01/2026 – 06/2026): M1.1 Špecifikácia požiadaviek; M1.2 Architektonický návrh; M1.3 GDPR DPIA draft.
\item WP2 (04/2026 – 12/2026): M2.1 API v1; M2.2 AI modul v1; M2.3 Sync v1.
\item WP3 (07/2026 – 06/2027): M3.1 UI prototyp; M3.2 Beta UI; M3.3 UX testy v1.
\item WP4 (07/2027 – 12/2028): M4.1 Pilot v1; M4.2 Report metrik v1; M4.3 Záverečná evaluácia.
\item WP5 (01/2026 – 12/2028): M5.1 Web \& dokumentácia; M5.2 Workshop(y); M5.3 Plán udržateľnosti.
\end{itemize}

% ===================== C =====================
\section*{VV–C: Rozpočet projektu (EUR)}
{\small
\begin{tabularx}{\textwidth}{@{}X>{\raggedleft\arraybackslash}p{2.2cm}>{\raggedleft\arraybackslash}p{2.2cm}>{\raggedleft\arraybackslash}p{2.2cm}>{\raggedleft\arraybackslash}p{2.6cm}@{}}
\rowcolor{apvvgray}\textbf{Položka} & \textbf{2026} & \textbf{2027} & \textbf{2028} & \textbf{Spolu}\\ \toprule
Priame prevádzkové náklady & 6~000 & 6~000 & 6~000 & 18~000\\
Mzdové a osobné náklady & 24~000 & 36~000 & 36~000 & 96~000\\
Sociálne a zdravotné odvody & 8~400 & 12~600 & 12~600 & 33~600\\
Cestovné náklady & 3~000 & 3~000 & 3~000 & 9~000\\
Materiál & 5~000 & 3~000 & 3~000 & 11~000\\
Služby (cloud, AI, bezpečnostný audit, UX) & 16~000 & 20~000 & 22~000 & 58~000\\
Energia, voda, komunikácie & 2~000 & 2~000 & 2~000 & 6~000\\
Popularizačné náklady & 2~000 & 3~000 & 3~000 & 8~000\\
Nepriame náklady & 12~000 & 15~000 & 15~000 & 42~000\\ \midrule
\textbf{Spolu podľa rokov} & \textbf{78~400} & \textbf{100~600} & \textbf{102~600} & \textbf{281~600}\\
Z toho: APVV / Spolufinancovanie & APVV: 239~360 & Cofin: 42~240 & — & —\\ \bottomrule
\end{tabularx}
}
% ===================== D =====================
\section*{VV–D: Harmonogram projektu a zoznam prístrojov}
{\small\begin{longtable}{@{}p{2.2cm}p{2.2cm}p{11.6cm}@{}}
\rowcolor{apvvgray}\textbf{Začiatok} & \textbf{Koniec} & \textbf{Názov etapy / míľniky}\\ \toprule
01/2026 & 06/2026 & WP1 – Analýza a návrh; M1.1 Špecifikácia, M1.2 Architektúra, M1.3 DPIA draft\\
04/2026 & 12/2026 & WP2 – Backend \& API; M2.1 API v1, M2.2 AI v1, M2.3 Sync v1\\
07/2026 & 06/2027 & WP3 – Frontend \& UX; M3.1 UI prototyp, M3.2 Beta UI, M3.3 UX testy v1\\
07/2027 & 12/2028 & WP4 – Pilot a evaluácia; M4.1 Pilot v1, M4.2 Report, M4.3 Záverečná evaluácia\\
01/2026 & 12/2028 & WP5 – Disseminácia a udržateľnosť; M5.1 Web, M5.2 Workshopy, M5.3 Udržateľnosť\\
\end{longtable}}

\noindent\textbf{Zoznam prístrojov v projekte}

\begin{tabularx}{\textwidth}{@{}lY@{}}
\Field{Vývojový server / cloud kredit}{CI/CD, staging, databázové prostredie, monitoring}
\Field{Testovacie mobilné zariadenia}{Android/iOS zariadenia pre overenie UX a notifikácií}
\end{tabularx}

% ===================== E =====================
\section*{VV–E: Čestné vyhlásenia}
Vyplní a podpíše štatutárny zástupca žiadateľa (ponechané na doplnenie pri podaní).

\vfill
\noindent\textit{Pozn.: Obsah vychádza z článku o projekte TIMA a z oficiálneho členenia formulára VV–A až VV–E.}
\end{document}
